\DocumentMetadata{
  pdfversion=2.0,
  pdfstandard=ua-2,
  lang=en-US,
  tagging=on,
  tagging-setup={math/setup=mathml-SE,tabsorder=structure}
}
\documentclass[11pt,letterpaper]{article}
\usepackage[margin=0.68in,includefoot]{geometry}
\usepackage{newcomputermodern}
\usepackage{amsmath,amscd,mathtools}
\providecommand{\square}{\mdlgwhtsquare}
\usepackage[hidelinks]{hyperref}
\hypersetup{
  pdftitle={Algebra PhD Exam, May 2023},
  pdfauthor={University of Florida Department of Mathematics},
  pdfsubject={Graduate mathematics examination},
  pdfdisplaydoctitle=true,
  bookmarksopen=true
}
\setcounter{secnumdepth}{0}
\setlength{\parindent}{0pt}
\setlength{\parskip}{0.28em}
\setlength{\emergencystretch}{3em}
\setlength{\leftmargini}{2.15em}
\setlength{\leftmarginii}{2.65em}
\setlength{\leftmarginiii}{3em}
\renewcommand{\labelenumi}{\textbf{\theenumi.}}
\pagestyle{empty}
\raggedbottom
\allowdisplaybreaks
\newenvironment{examproblems}[1]{%
  \begin{enumerate}%
  \setcounter{enumi}{#1}%
  \setlength{\topsep}{0.35em}%
  \setlength{\partopsep}{0pt}%
  \setlength{\itemsep}{0.48em}%
  \setlength{\parsep}{0.18em}%
}{\end{enumerate}}
\newenvironment{examsubparts}{%
  \begin{enumerate}%
  \setlength{\topsep}{0.18em}%
  \setlength{\partopsep}{0pt}%
  \setlength{\itemsep}{0.16em}%
  \setlength{\parsep}{0.1em}%
}{\end{enumerate}}
\newenvironment{examinstructions}{%
  \par\begingroup\itshape
}{%
  \par\endgroup\vspace{0.18em}
}
\makeatletter
\renewcommand\section{\@startsection{section}{1}{0pt}%
  {-1.25ex plus -.25ex minus -.1ex}%
  {0.55ex plus .1ex}%
  {\normalfont\large\bfseries}}
\renewcommand{\@maketitle}{%
  \newpage\null\vskip -1.2em%
  {\footnotesize\bfseries
    \href{https://www.ufl.edu/}{University of Florida}, \href{https://math.ufl.edu/}{Department of Mathematics}\par}%
  \vskip 0.18em%
  \tagstructbegin{tag=Title}%
    {\LARGE\bfseries\href{https://gma.math.ufl.edu/exams/algebra-phd/}{Algebra PhD Exam}, May 2023\par}%
  \tagstructend%
  \vskip 0.35em%
  \tagmcbegin{artifact}%
    \hrule height 1pt%
  \tagmcend%
  \vskip 0.55em%
}
\makeatother
\title{Algebra PhD Exam, May 2023}
\author{}
\date{}
\begin{document}
\maketitle
\thispagestyle{empty}
\begin{examinstructions}
Answer seven problems. Write your answers clearly in complete English sentences. You may quote results (within reason) as long as you state them clearly.
\end{examinstructions}
\begin{examproblems}{0}
\item
Let~\(K\) be the splitting field of \(x^7-2\) over~\(\mathbb{Q}\). What is \([K:\mathbb{Q}]\)? Justify your answer.
\item
Give an example of a Galois extension of \(\mathbb{Q}(\zeta_5)\) with degree~\(5\) (where \(\zeta_5\) is a primitive 5th root of unity). Give an example of a Galois extension of~\(\mathbb{Q}\) with degree~\(5\). Justify your answers.
\item
Let~\(K\) be a field and~\(G\) a finite subgroup of \(\operatorname{Aut}(K)\).
\begin{examsubparts}
\item[\textnormal{(i)}]
What does the “theorem of the fixed field” say about \([K:K^G]\)?
\item[\textnormal{(ii)}]
Use the theorem of the fixed field to show that if \(K/F\) is a finite extension of fields then \(\#\operatorname{Aut}(K/F)\leq [K:F]\) with equality if and only if~\(F\) is the fixed field of \(\operatorname{Aut}(K/F)\).
\item[\textnormal{(iii)}]
Use the theorem of the fixed field to show that if \(F=K^G\) then \(K/F\) is a Galois extension with Galois group~\(G\).
\end{examsubparts}
\item
Let \(\mathrm{Groups}\) denote the category of groups and \(\mathrm{Sets}\) denote the category of sets.
\begin{examsubparts}
\item[\textnormal{(i)}]
Use free groups to construct a left adjoint to the forgetful functor \(U:\mathrm{Groups}\to\mathrm{Sets}\). (You can skip checking “naturality” i.e. compatibility with morphisms.)
\item[\textnormal{(ii)}]
Show that the functor \(F:\mathrm{Groups}\to\mathrm{Sets}\) defined by 
\[
F(G)=\{(x,y)\in G\times G:xyx^{-1}y^{-1}=1\}
\]
 is a representable functor. (Again, you can skip checking “naturality”.)
\end{examsubparts}
\item
Let~\(k\) be a field. For a~\(k\)-vector space~\(V\), recall that the dual vector space \(V^*\) is defined to be the~\(k\)-vector space of linear functionals on~\(V\), i.e. linear transformations from~\(V\) to~\(k\).
\begin{examsubparts}
\item[\textnormal{(i)}]
Construct a natural linear transformation \(\psi:V^*\otimes_k W^*\to(V\otimes_k W)^*\).
\item[\textnormal{(ii)}]
Show that~\(\psi\) is an isomorphism if~\(V\) and~\(W\) are finite dimensional.
\end{examsubparts}
\item
Let~\(R\) be a unital ring.
\begin{examsubparts}
\item[\textnormal{(i)}]
State at least two equivalent definitions of a projective~\(R\)-module. (You don’t need to show they are equivalent, but you should include any definition you want to use in the next part.)
\item[\textnormal{(ii)}]
Let~\(M\) and~\(N\) be~\(R\)-modules with \(M\subset N\). If~\(M\) is projective and \(N/M\) is projective, prove that~\(N\) is projective.
\end{examsubparts}
\item
Are the functors \(F(A)=\operatorname{Hom}_{\mathbb{Z}}(A,\mathbb{Q}/\mathbb{Z})\) and \(G(A)=\operatorname{Hom}_{\mathbb{Z}}(A,\mathbb{Z})\) from the category of~\(\mathbb{Z}\)-modules to the category of~\(\mathbb{Z}\)-modules exact? Justify your answer.
\item
Prove that the intersection of all prime ideals in a commutative ring is the nilradical. (Suggestion: for one direction, given a non-nilpotent element~\(a\) consider the set of all ideals not containing any power of~\(a\).)
\item
Let~\(\mathbb{R}\) and~\(\mathbb{C}\) denote the real and complex numbers as usual.
\begin{examsubparts}
\item[\textnormal{(i)}]
Give an example of two different maximal ideals of \(\mathbb{R}[x]\) which have the same zero set in the affine line \(\mathbb{A}^1_{\mathbb{R}}\).
\item[\textnormal{(ii)}]
Give an example of two ideals of \(\mathbb{C}[x]\) which have the same zero set in \(\mathbb{A}^1_{\mathbb{C}}\).
\item[\textnormal{(iii)}]
State Hilbert’s Nullstellensatz.
\end{examsubparts}
\item
Let~\(R\) be a commutative ring and~\(D\) a multiplicatively closed subset of~\(R\) containing~\(1\). State the universal property for the localization \(D^{-1}R\), show that any two rings satisfying the universal property are isomorphic, and construct a ring satisfying the universal property. (Checking the construction is well-defined will involve a variety of verifications which you can skip. But do check it satisfies the universal property.)
\item
Give an example of the following, and briefly justify your answers:
\begin{examsubparts}
\item[\textnormal{(i)}]
A discrete valuation ring whose field of fractions is not of characteristic zero.
\item[\textnormal{(ii)}]
A Dedekind domain that is not a unique factorization domain.
\end{examsubparts}
\end{examproblems}
\end{document}
